\section{Failure of compactness and effective completeness}
\label{sec:second-order-logic-04}

\begin{theorem}[Failure of compactness under full semantics]
Full second-order semantic consequence is not compact.
\end{theorem}

\begin{proof}
Expand \(\Lang_{\mathsf N}\) by a constant \(c\).  Let \(T\) contain the full
second-order Peano axioms and, for every \(n\in\N\), the sentence
\[
  c\ne S^n0.
\]
Every finite subset of \(T\) is satisfied under full semantics by a suitable
structure: use the standard natural numbers and interpret \(c\) as a number
larger than every index occurring in that finite subset.  If one structure
satisfied all of \(T\) under full semantics, categoricity would identify its
domain with \(\N\), while the displayed sentences would require \(c\) to
differ from every natural number.  Hence no such structure exists.
\end{proof}

\begin{theorem}[No effective complete calculus for full validity]
Here an effective calculus means one whose finite derivations are
computably checkable, so its theorems form a computably enumerable set.  There
is no such calculus \(\mathsf D\) for which, for every second-order sentence
\(\sigma\),
\[
  \mathsf D\derives\sigma
  \quad\Longleftrightarrow\quad
  \entails_{\mathrm{full}}\sigma.
\]
Indeed, the set of full second-order validities is not computably enumerable.
\end{theorem}

\paragraph{Proof.}
The complete proof, including the effectiveness of the reduction and the
equivalence below, is \cref{thm:p02-w02-p02-040-no-effective-calculus}.
Use the finite full second-order axiomatization of arithmetic in the signature
with \(0,S,+,\times\).  For each first-order arithmetic sentence \(\sigma\),
effectively form the second-order sentence
\(\mathsf{PA}_2\to\sigma\).  Full categoricity gives
\[
  \entails_{\mathrm{full}}(\mathsf{PA}_2\to\sigma)
  \quad\Longleftrightarrow\quad
  (\N,0,S,+,\times)\entails\sigma.
\]
If full second-order validities were computably enumerable, true
first-order arithmetic would therefore be computably enumerable,
contradicting the corollary of Tarski's undefinability theorem
(\cref{thm:tarski-undefinability}) that true arithmetic is not computably
enumerable.  The reduction, its effectiveness and the computability step
through \cref{thm:p04-ce-sigma-one} are proved in
\cref{thm:p02-w02-p02-040-no-effective-calculus}.

This theorem excludes an effective, sound, and complete calculus for full
validity.  It does not exclude an ineffective presentation that takes all full
validities as axioms.  The effectiveness qualification is therefore part of
the theorem, not an optional refinement; see
\cite{BarwiseFeferman1985,Shapiro1991}.

\begin{theorem}[Categoricity of the full second-order reals]
\label{thm:second-order-reals-categorical}
Let \(\Lang_{\mathsf R}=\set{0,1,+,\times,<}\).  Take the finitely many
first-order axioms stating that \((0,1,+,\times,<)\) is an ordered field,
together with the second-order completeness sentence
\[
  \forall X\,
  \Bigl[
    \exists x\,X(x)
    \land
    \exists y\,\forall x\,(X(x)\to x\le y)
    \;\to\;
    \exists z\,
    \bigl(
      \forall x\,(X(x)\to x\le z)
      \land
      \forall y\,(\forall x\,(X(x)\to x\le y)\to z\le y)
    \bigr)
  \Bigr],
\]
where \(u\le v\) abbreviates \(u<v\lor u=v\) in this scope: every nonempty
subset with an upper bound has a least upper bound.  Any two
\(\Lang_{\mathsf R}\)-structures satisfying these axioms under full
semantics are isomorphic.  In particular, every such structure is isomorphic
to the ordered field \(\R\), and its domain is uncountable.
\textit{Status: proved in \cref{thm:p02-w02-p02-040-reals-categorical}.}
\end{theorem}

The ordered-field axioms abbreviated above are listed in
\cref{def:p02-w02-p02-040-ordered-field}.  The categoricity proof, which sends
each element of one model to the element of the other determined by the same
cut of the common rationals and verifies that this map is an ordered-field
isomorphism, and the proof of uncountability are given in
\cref{thm:p02-w02-p02-040-reals-categorical}; compare \cite{Shapiro1991}.

Full second-order logic also lacks the first-order L\"owenheim--Skolem
properties.  The conjunction of the displayed second-order Peano axioms is
satisfied by no uncountable structure
under full semantics, so upward L\"owenheim--Skolem fails.  In the other
direction, the axioms of \cref{thm:second-order-reals-categorical} are
satisfied by the infinite structure \(\R\), while
\cref{thm:second-order-reals-categorical} makes every structure satisfying
them under full semantics isomorphic to \(\R\) and therefore uncountable.
These axioms are thus satisfied by an infinite structure but by no countable
structure; downward L\"owenheim--Skolem fails.

\begin{theorem}[Lindstr\"om]
\label{thm:lindstrom}
Let \(\mathcal L\) be a regular abstract logic extending first-order
logic: one closed under the Boolean connectives, quantification, and
relativization, with satisfaction invariant under isomorphism.  If
\(\mathcal L\) satisfies countable compactness and the downward
L\"owenheim--Skolem property to countable structures, then
\(\mathcal L\) has exactly the expressive power of first-order logic.
\textit{Status: cited theorem.}
\end{theorem}

See \cite{Lindstrom1969} for the original proof and
\cite[Chapter~XIII]{EbbinghausFlumThomas1994} for a complete textbook
treatment, including the exact closure conditions abbreviated here as
regularity.  The theorem explains the pattern of this section: full
second-order logic strictly exceeds first-order expressive power, so it
must lose countable compactness or the downward L\"owenheim--Skolem
property, and this section shows that it loses both.

